\begin{frame}{Why timing matters in TOF-PET}
\small

\begin{columns}[T,totalwidth=\textwidth]

\begin{column}{0.48\textwidth}

\begin{block}{TOF-PET imaging}
A positron annihilation produces two photons emitted in opposite directions.
Two detectors measure the corresponding signals and estimate the
\textbf{difference in photon arrival time (TOF)}.
\end{block}

\vspace{0.25em}

\begin{alertblock}{Why better timing helps}
The arrival-time difference constrains the annihilation position along the line connecting the two detectors:

\[
\Delta x = \frac{c\,\Delta t}{2}.
\]

\vspace{0.15em}
A more precise estimate of $\Delta t$ means a more precise localization of the event.
\end{alertblock}

\end{column}


\begin{column}{0.48\textwidth}
\centering

\smartimage[0.94\linewidth]{0.57\textheight}
{figures/TOF.png}
{TOF-PET principle: the difference in photon arrival times localizes the annihilation event along the line of response.}

\vspace{0.3em}

\begin{block}{Project question}
\centering
What is the best timing performance achievable,
and how can we approach it with a practical and scalable method?
\end{block}

\end{column}

\end{columns}

\end{frame}
